\documentclass[a4paper,10pt]{report}
\usepackage{fancyhdr}
\usepackage{fancybox}
\usepackage{tipa}
\usepackage{framed}
\usepackage[retainorgcmds]{IEEEtrantools}
\usepackage[utf8x]{inputenc}
\usepackage{epsfig}
\usepackage{anysize}
\usepackage{gensymb}
\usepackage{amssymb}
\usepackage{graphicx}
\usepackage{caption}
\usepackage{subcaption}
\usepackage[utf8x]{inputenc}
\usepackage{longtable}
\usepackage{url}
\usepackage{plain}
\renewcommand{\bibname}{Refrences}

\usepackage[plainpages=false]{hyperref}
\hypersetup{
	colorlinks=true,
	linkcolor=black,
	citecolor=black,
	filecolor=black,
	urlcolor=black,
}
\lhead{}

\setcounter{page}{-100}
\setcounter{secnumdepth}{5}

\newenvironment{fminipage}%
{\begin{Sbox}\begin{minipage}}%
{\end{minipage}\end{Sbox}\fbox{\TheSbox}}
\title{\textbf{Secure Document Sharing using RWFM}}
\vspace{2.5cm}
\author{\bf{A R \& D Project Report}\\
		\\
		\emph{Submitted in partial fulfillment of requirements for the degree of}\\
		\\
		\bf{Master of Technology}\\
		\\
		\emph{by}\\
		\\
		    \bf{Asif Ali}\\
		\bf{Roll No : 143059009}\\
		\\
		\emph{under the guidance of}\\
		\\
		\bf{Prof. R.K. Shyamsundar}\\
		\\\\\
		\includegraphics[height=3.5cm]{./images/iitb_logo}\\
		\\
		\bf{Department of Computer Science and Engineering}\\
		\bf{Indian Institute of Technology,Bombay}\\}
\date{}
\marginsize{2.5cm}{2.5cm}{1cm}{1.5cm}
\setlength{\headheight}{15pt}
\pagestyle{fancy}
\usepackage{graphicx}
\begin{document}
\maketitle
\pagenumbering{roman}
\tableofcontents
\begin{abstract}
Most of our world are becoming online like social intractions, transferrence of thoughts
and ideas takes place through internet. These ideas are shared using documents using document files,
videos,music files etc. On transferring these files, security of these transfer became very
important aspect. Most of applications follows some rudimentary methods to secure these transferrence
but implementation of these methods are very complex.So, it will be difficult to manipulate
the already existing policy.On the other hand, there is a method called RWFM which is simple to design
and update security policies.
\end{abstract}
\pagenumbering{arabic}
\chapter{Introduction}

\label{chapter:intro}
\end{document}
